\section{Labors of Hercules, Part 4}

The Erymanthian Boar\footnote{\url{https://en.wikipedia.org/wiki/Erymanthian\_Boar}}

Hercules had been given the (originally) ten labors as a penance for letting Hera drive him to wrath, in which he slew his six sons. So don't let anyone tell you that penance is a ``Catholic'' thing. All of the labors center around a primordial world of monsters that will be rededicated to the sun-God Apollo, or to Zeus, from whom the power to cleanse \& purify this region of the world comes.

Hercules originally received the wrath of Hera because he was the illegitimate son of Zeus and Alcmene. He is half-born, a godling, immortal by birth right, but not in actuality. Hera is one of the daughters of Chronos, but also a wife to Zeus. Hera represents an older form of matriarchal worship connected with earth rites and the underworld (like Demeter) that is incorporated into the male \& solar spirit of the ``new gods'' of Greece. Her typical function is rage against Zeus' consorts and love affairs. We can view Hera as a suppressed kind of female power that tends to take vengeance on the half-god offspring of Zeus' spirit mingling (like the wind) with mere mortals. Not only is he creating a race that is superior to those who worship Hera, the warrior and martial ethos threatens the earlier balance. As such, Hera is ambiguous. She is both a reminder that the Greeks ``breathed their truth through lies'' in the new myths (a fact Clement of Alexandria alludes to), and also a principle of retardation: there is no ``going back'' to the earlier stage.

We have to view this evolution from a solar perspective. On the one hand, the revolt of the warriors against the priests heralded a new Dark Age. On the other, man never rises except by suffering. As the ancients would say ``he suffers into Truth''. Hera is fighting a lost cause. The forces of chaos have penetrated the ancient ordering of the castes, turning each against the others. Corruption begins with the priests, who compromise (out of fear) with the warrior caste, which is ascendent.
\emph{The priests should suffer like Prometheus to bring them into a proper order.} They lash out with curses and penances, after the betrayal of their compromise, which they should never have agreed to (we see the underground Occult engaged in such a struggle today). An analogous event occured in the West, in which the Frankish nobility went over to the Revolution through decadence and compromise, with the priests like Abbe Sieyes courting natural nobles such as Mirabeau, who of course caters to the third estate, who will later flirt with the masses. Garibaldi and his ilk (natural \emph{aristoi}) become leaders of ``Revolution''. \emph{As Cologero has written, we can see in these modern events Time greatly sped up, during the Kali Yuga, for our own edification, if we could but read aright. It is no accident that the natural dramatic flair of the French was interested in Classical themes!}

Some men, like Hercules, will rise to the challenge, and prove themselves, redeeming the turbulence of the purely horizontal swirl of the world-snake Ourobous into a vertical ascent. Evolution always operates: the snake writhes and awaits the champion to bend it to a will. Even today, our genetic structures are engaged in such programming\footnote{\url{http://www.koanicsoul.com/blog/2013/06/16/edward-snowden-a-true-neanderthal-shivs-the-melonhead-nwo/}}; they await men capable of rising, like eagles on a vertical lift of heated wind, up into the heavens. This ascension has to involve Christian ``ideas'' (not necessarily formally) \emph{as it is only through sacrifice of some sort that the necessary purity is invoked}. \textbf{Even in Hercules' time, penance is necessary.} That is why Hera is still important -- the solar spirit embraces Fate, suffers into Truth, and guides it upward. The true man sees in the Modern World an immense potential opportunity, and begins the ascent, without for all that, endorsing the degeneracy he finds around him. Hera represents the Karma of past sins, in this case, ``sins of the Father''.

Hera is ``the adversary'': not an enemy per se, but someone (who for the strong man) is fulfilling a purpose \& role. Without the consent of Hera, Hercules cannot be deified (as we shall she). It is her job to test the hero, and although she is insufficient to initiate solar truth, she has the power to test it.

The great boar is tracked through the snow, driven over the mountain ranges by Hercules, who descends upon it relentlessly, driving it before him. Chiron the centaur had advised him to track it in the snow. Hercules is now succeeding more easily in his labors, and the boar proves ``easy'' to track down. The hero is gaining more power.

What is interesting is the tale of the centaurs. While visiting Pholus during the hunt, Hercules desires wine, and convinces Pholus to open the one vessel he has, a gift from Dionysius. The centaurs do not know how to drink moderately, and end up becoming drunk and attacking Hercules, who drives them off with arrows. Chiron is hit, and the pain it causes this immortal makes him volunteer to replace Prometheus on the mountain, being tortured by the eagle. Hercules shoots the eagle, to deliver Chiron. This teaches us that Hercules is receiving the immortality which Chiron was unable to utilize properly. The drunken man-beast who tragically frees Prometheus is imbuing the rising hero with his powers.

Hercules has already conquered the passions, so he is worthy to drink of Dionysius' wine, and can easily disperse the dangers arising thereby. His arrows, dipped in the Hydra's blood, are more than equal to chastening the centaurs, who are wise up to a point, but unruly.

There is only so much immortality floating around at one point of time. Those who misuse this gift (their talents) will have it given to someone else more worthy. This should be an encouragement to everyone who is a spiritual seeker. When you see the powerful and ignorant, or the mighty and the evil, squandering their gifts, you should know that these will become available to you, if they continue to defy the deathless and immortal source of these divine blessings. ``The meek inherit the earth''. This is taught explicitly in the Psalms, and endorsed in the Beatitudes and apostolic teaching. ``Run the race, as if to win''. When someone releases their anger unjustly upon you, a right response will not only help them, it will put their lost energy at your disposal. Behind the ``meekness'' of the Christian tradition is hidden an ancient teaching concerning the accumulation of power. When Saint Peter crosses himself in the presence of the emperor Nero, causing Simon Magus to fall from the sky (Magus was demonstrating the power of his magic), the lost energy will go to Saint Peter. This happens with Cyprian the Mage \& Justiana\footnote{\url{https://en.wikipedia.org/wiki/Cyprian\_and\_Justina}}, as well (cited by Tomberg in \emph{Meditations}).

Since we are discussing magic at Gornahoor, it would be helpful to meditate on how to make it as powerful as possible: counter-intuitively, it would seem that purity is that which renders it deeper and whiter magic. This purity takes different guises, but purity it is, nonetheless, \& necessary to avoid the trap indicated by Dante, who places Simon Magus in hell. ``To whom much is given, much is required.''

The tale ends humorously: the Greeks loved to depict the scene in which Hercules presents the boar to the tyrant, only to have him hide inside of a large vase. These endings are further vindication of Hercules' divine authority and power, and an important ``unmasking'' of the powers-that-be: they are unequal to even receive the gift of Hercules, which they asked for.

The legend of Hercules remains a beautiful re-telling of what was good in the pre-Olympian world, a transmuting of it through the might of Hercules into something even higher and even better. Something similar occurred in the West, and one can discern the operations of Christianity on the classical ideals in much the same manner. In each case, the essence or import is preserved, but the eagle makes another turn in the gyre of the circle, and achieves greater altitude. Like the hero, it rises, on steady wings and a strong upward draft. For an instance of this, see Matthew Arnold's essay on Maurice du Guerin's tale of the centaurs\footnote{\url{https://books.google.com/books?id=BOY-AAAAIAAJ\&pg=PA64\&lpg=PA64\&dq=centaurs+matthew+arnold+guerin\&source=bl\&ots=y2r8ffCQmx\&sig=xp9Sua44p9O5Pkqht9IeYGycVW8\&hl=en\&sa=X\&ei=gqcgUpmnKam\_sATF24GgCg\&ved=0CCwQ6AEwAA\#v=onepage\&q=centaurs\%20matthew\%20arnold\%20guerin\&f=false}}.

Hercules represents the new warrior, the best of that breed, who dare to will a union of warrior-priest, in order to achieve immortality. One might say that a warrior is potentially an aspirant priest: if he succeeds, if he navigates the divine-demonic world successfully, he discharges his penance (against which and with which he fights), achieving the essence of immortality. Later on, in the West, this baptized chivalric ideal would be made more explicit in the body of the king\footnote{\url{http://www.libraryofsocialscience.com/ideologies/docs/the-kings-two-bodies/index.html}}. The loss of the primordial world, its descent into chaos, has to be \emph{redeemed}, through the storming of heaven. If Christianity had not existed, it would have had to be invented. For those willing to undertake a similar labor, modern Christianity has its primordial monsters and jealous and offended goddesses which are willing to oblige you, in standing the test.


\flrightit{Posted on 2013-08-30 by Logres}

\begin{center}* * *\end{center}

\begin{footnotesize}
\begin{sffamily}
\texttt{William on 2013-09-30 at 08:08 said:}

``When someone releases their anger unjustly upon you, a right response will not only help them, it will put their lost energy at your disposal.'' Is there anyplace where one can read further about the idea of transmuting negative energy directed at oneself? (As opposed to simply withstanding it.) Thanx!

\hfill

\texttt{scardanelli on 2013-09-30 at 09:33 said:}

William,

Mouravieff discusses this in his Gnosis series, and I assume that Ouspensky and Gurdjieff do as well, though I have never read the later. I think it's mentioned in the third Gnosis book, though I don't have them in front of me. Logres can probably go into more detail, but from my understanding, when one remains conscious and aware and doesn't react to negative emotions but stays firmly within the realm of silence (or magnetic center in Mouravieff's terms), one doesn't expend energy in mechanical reaction but absorbs the other's output of energy in a feeling of freedom and love. It has been a little while since I read Mouravieff though, so perhaps others could correct any of my misunderstandings.


\hfill

\end{sffamily}
\end{footnotesize}

